\begin{theorem}[Multiplicity-one semi-regular representation]
    \label{thm:peps_regular_minimal_representation}
    \lean{TNLean.PEPS.exists_minimalSemiRegular_leftRegular_fourier}
    \leanok
    \uses{thm:peps_regular_unitary_fourier,
        thm:peps_semiregular_group_algebra,
        thm:peps_torus_block_multiplicity_state}
    For every finite group, the Fourier representations in
    (\ref{eq:peps_regular_unitary_fourier}) can be chosen so that
    $U_{\min}(g)=\bigoplus_i D_i(g)$ is unitary and semi-regular.
    The multiplicity of each character $\chi_i$ in $U_{\min}$ is one.
    The same chosen matrices give the actual unitary Fourier decomposition
    $U_{\mathrm{reg}}\cong\bigoplus_i D_i\otimes I_{d_i}$.
    Thus the group determines the multiplicity-one representation used
    in \cite[Section~7]{Schuch2010PEPS}.
\end{theorem}
\begin{proof}\leanok
    The direct sum is unitary because each $D_i$ is unitary. If a linear
    combination of the minimal group matrices vanishes, it vanishes in
    every block $D_i$. Repeating each zero block $d_i$ times and changing
    back by the actual Fourier basis gives the same zero combination of
    regular group matrices. Those matrices are linearly independent.
    Thus the minimal matrices are linearly independent and their
    group-algebra action is faithful, which implies semi-regularity.
    Finally, its character is $\sum_i\chi_i$. Character orthogonality
    and pairwise inequivalence give multiplicity one for each $\chi_i$.
\end{proof}
